% Draft English CV — edit locally and recompile with:
%   ./build_cv.sh
% or: xelatex Masato_Tsutsumi_CV.tex && cp Masato_Tsutsumi_CV.pdf ../static/uploads/
\documentclass[10pt,a4paper]{article}

\usepackage[margin=1.35cm]{geometry}
\usepackage{fontspec}
\usepackage{xcolor}
\usepackage{titlesec}
\usepackage{enumitem}
\usepackage{tabularx}
\usepackage{hyperref}
\usepackage{microtype}
\usepackage{parskip}

\setmainfont{Helvetica Neue}[
  UprightFont = * ,
  BoldFont = * Bold ,
  ItalicFont = * Italic ,
  BoldItalicFont = * Bold Italic
]

\definecolor{ink}{HTML}{0F172A}
\definecolor{accent}{HTML}{1E3A5F}
\definecolor{muted}{HTML}{475569}
\definecolor{rule}{HTML}{CBD5E1}

\hypersetup{
  colorlinks=true,
  urlcolor=accent,
  pdfauthor={Masato Tsutsumi},
  pdftitle={Curriculum Vitae — Masato Tsutsumi, Ph.D.}
}

\pagestyle{empty}
\setlength{\parindent}{0pt}
\setlength{\parskip}{0.18em}
\setlist[itemize]{leftmargin=1.05em, itemsep=0.08em, topsep=0.12em, parsep=0pt}
\setlist[enumerate]{leftmargin=1.2em, itemsep=0.12em, topsep=0.15em, parsep=0pt}

\titleformat{\section}
  {\normalsize\bfseries\color{accent}}
  {}{0pt}{\MakeUppercase}[\vspace{-0.28em}{\color{rule}\titlerule[0.7pt]}\vspace{0.2em}]
\titlespacing*{\section}{0pt}{0.55em}{0.18em}

\newcommand{\cvitem}[2]{%
  \noindent\begin{tabularx}{\textwidth}{@{}>{\raggedright\arraybackslash}p{2.7cm} X@{}}
    \textcolor{muted}{\footnotesize #1} & \small #2 \\
  \end{tabularx}\vspace{0.04em}
}

\newcommand{\talk}[2]{%
  \noindent\begin{tabularx}{\textwidth}{@{}>{\raggedright\arraybackslash}p{2.35cm} X@{}}
    \textcolor{muted}{\footnotesize #1} & \footnotesize #2 \\
  \end{tabularx}\vspace{0.02em}
}

\begin{document}

\noindent
\begin{minipage}[t]{0.66\textwidth}
  {\fontsize{20}{24}\selectfont\bfseries\color{ink} Masato TSUTSUMI, Ph.D.}\\[0.25em]
  {\large\color{accent} Assistant Professor}\\[0.15em]
  {\small\color{muted}
  Life Science Center for Survival Dynamics (TARA)\\
  University of Tsukuba, Japan}
\end{minipage}%
\hfill
\begin{minipage}[t]{0.32\textwidth}
  \raggedleft\footnotesize\color{muted}
  \href{mailto:mtsutsumi@tara.tsukuba.ac.jp}{mtsutsumi@tara.tsukuba.ac.jp}\\
  \href{https://masa10223.github.io/en/}{masa10223.github.io}\\
  \href{https://scholar.google.co.jp/citations?user=sLCyYs8AAAAJ}{Scholar} $\cdot$
  \href{https://github.com/masa10223}{GitHub} $\cdot$
  \href{https://researchmap.jp/masato_tsutsumi?lang=en}{researchmap}\\
  Updated: August 2026
\end{minipage}

\vspace{0.4em}
{\color{accent}\hrule height 1.8pt}
\vspace{0.35em}

{\small
I study life phenomena as a hierarchical system of morphology, psychology, and behavior, and quantify them as information. My work combines deep generative models for morphological quantification, mathematical models of psychological states (e.g., social stress), and quantitative analysis of inter-individual behavior, aiming to uncover universal principles that describe biological state and dynamics.
}

\section{Appointments}
\cvitem{2025/12 -- present}{\textbf{Assistant Professor}, Life Science Center for Survival Dynamics (TARA), University of Tsukuba\\
\emph{Mathematical Life Dynamics Project} (Assoc.\ Prof.\ Nen Saito)}
\cvitem{2025/12 -- present}{\textbf{Visiting Researcher}, Graduate School of Medicine, Nagoya University\\
\emph{Data-Driven Biology Laboratory} (Prof.\ Naoki Honda)}
\cvitem{2024/10 -- 2025/11}{\textbf{Designated Assistant Professor}, Graduate School of Medicine, Nagoya University}
\cvitem{2023/10 -- 2024/09}{\textbf{Designated Researcher}, Graduate School of Integrated Sciences for Life, Hiroshima University}
\cvitem{2020/04 -- 2023/03}{\textbf{Graduate Student Research Associate}, RIKEN}

\section{Education}
\cvitem{2020/04 -- 2023/09}{\textbf{Ph.D.\ in Science}, Graduate School of Science, The University of Tokyo (Supervisor: Prof.\ Chikara Furusawa)}
\cvitem{2018/04 -- 2020/03}{\textbf{M.S.\ in Science}, Graduate School of Science, The University of Tokyo (Supervisor: Prof.\ Chikara Furusawa)}
\cvitem{2014/04 -- 2018/03}{\textbf{B.S.\ in Science}, Department of Physics, Faculty of Science, The University of Tokyo}

\section{Teaching}
\cvitem{2025/06 -- 2025/07}{\textbf{Medical Data Science I}, Nagoya University School of Medicine}
\cvitem{2024/10 -- 2025/03}{\textbf{Basic Medical Seminar}, Nagoya University School of Medicine}
\cvitem{2019/04 -- 2019/08}{\textbf{Computational Experiments I} (Teaching Assistant), Department of Physics, The University of Tokyo}

\section{Grants \& Awards}
\cvitem{2025}{\textbf{Takeda Science Foundation Medical Research Grant} (Basic Research)}

\section{Publications}
\begin{enumerate}
  \item Yamaguchi, M., \textbf{Tsutsumi, M.}, Kuroda, Y., Soeda, Y., Yamaguchi, T. \textit{Artificial-Intelligence-Based Cephalometric Landmark Detection in Lateral Cephalograms}. \textit{Journal of Clinical Medicine} (2026). \href{https://doi.org/10.3390/jcm15155943}{doi:10.3390/jcm15155943}
  \item Uchida, Y., \textbf{Tsutsumi, M.}, Ichii, S., Irie, N., Furusawa, C. \textit{Deciphering the origin of developmental stability: The role of intracellular expression variability in evolutionary conservation}. \textit{Evolution \& Development} (2024).
  \item \textbf{Tsutsumi, M.}, Saito, N., Koyabu, D., Furusawa, C. \textit{A deep learning approach for morphological feature extraction based on variational auto-encoder: an application to mandible shape}. \textit{npj Systems Biology and Applications} (2023). \href{https://doi.org/10.1038/s41540-023-00293-6}{doi:10.1038/s41540-023-00293-6}
\end{enumerate}

\section{Preprints}
\begin{itemize}
  \item \textbf{Tsutsumi, M.}, Kubo, T., Kato, T.A., Honda, N. Mind the gap: quantifying population--individual gap in depressive symptom dynamics through energy landscapes. \textit{bioRxiv} (2026; v2).
  \item Matsuda-Watanabe, N., \textbf{Tsutsumi, M.}, Okumura, M., Chihara, T. Chemosensory input suppresses cannibalism by stabilizing social feeding boundaries in Drosophila larvae. In submission (2026).
  \item \textbf{Tsutsumi, M.}, Saito, N., Yamaguchi, T., Sasaki, T., Furusawa, C. Decoding Anadara shell morphology with deep learning. \textit{bioRxiv} (2026).
  \item Sakaguchi, S., \textbf{Tsutsumi, M.}, Nishi, K., Honda, N. Disentanglement of batch effects and biological signals across conditions in the single-cell transcriptome. \textit{bioRxiv} (2025).
\end{itemize}

\section{Invited Talks}
\talk{2026/07/30}{AI for measuring morphological differences among living organisms. \emph{FY2026 University of Tsukuba Open Lecture: Frontiers of Interdisciplinary Life Science Research 2026}.}
\talk{2026/07/27}{Towards quantification of form using information science. \emph{Workshop on Quantification of Form 2026}.}
\talk{2026/07/15}{Development of methods for quantifying biological morphology with an information-science approach. \emph{78th Annual Meeting of the Japan Society for Cell Biology}.}
\talk{2026/06/09}{Variational energy landscape analysis reveals population trends and individual heterogeneity during and after COVID-19 emergency declarations. \emph{40th Annual Conference of the Japanese Society for Artificial Intelligence (JSAI 2026)}.}
\talk{2026/03/05}{Towards quantification of diverse life phenomena: morphology, behavior, and psychology as information. \emph{SSTB2026 - Spring School for Theoretical Biology 2026}.}
\talk{2024/09/16}{Development of quantitative analysis methods for biological morphology using deep learning. \emph{Japan Society for Industrial and Applied Mathematics Annual Meeting 2024}.}
\talk{2024/03/09}{Development of quantitative analysis methods for biological morphology using deep learning. \emph{Online Seminar Series ``Future of Evolutionary Biology with AI''}.}
\talk{2024/02/14}{Development of quantitative analysis methods for biological morphology using deep learning. \emph{8th Theoretical Immunology Workshop}.}
\talk{2022/11/10}{Quantification of various biological morphologies using machine learning. \emph{Hiroshima University Mathematical Life Sciences Program Seminar}.}
\talk{2021/10}{Quantification and application of biological morphology using machine learning: Focusing on primate mandibles. \emph{National Institute for Basic Biology Departmental Seminar}.}

\section{Conference Presentations}
\talk{2026/07/12}{Mind the gap: quantifying population--individual differences in depressive symptom dynamics using energy landscapes. \emph{2nd Systems Behavioral Science Workshop}.}
\talk{2026/06/27}{Morphological character selection using machine learning for phylogenetic analysis: revisiting the phylogenetic position of lepidosaurs. \emph{Annual Meeting of the Palaeontological Society of Japan 2026}.}
\talk{2026/02/20}{Quantitative Analysis of Biological Morphology and Behavior Based on Deep Learning: Development of Morpho-VAE and DOLO. \emph{Frontiers of coevolutionary phenotypic emergence research: tracing how research unfolds from its roots}.}
\talk{2026/01/13}{Analysis of morphology and behavior using deep learning methods. \emph{Principles for Creating Life Functions through Evolutionary Information Assembly}.}
\talk{2025/12/05}{Development of a markerless multi-individual tracking system for controlling cannibalistic behavior in crickets. \emph{47th Annual Meeting of the Japanese Society for Comparative Physiology and Biochemistry}.}
\talk{2025/11/30}{Morphological quantification using deep learning: Morpho-VAE. \emph{99th Symposium on Science of Form ``Shape that brings motion''}.}
\talk{2025}{Visualization of morphological variation using 3D Variational Autoencoder. \emph{34th Annual Meeting of the Japanese Society for Bioimaging}.}
\talk{2025/05/30}{Behavioral analysis of double mutant Drosophila larvae using DOLO (Drosophila tracking with YOLO). \emph{1st Systems Behavioral Science Workshop}.}
\talk{2025/02/20}{Estimation of psychological state transitions during/after emergency declarations using energy landscape analysis. \emph{SSTB2025 - Spring School for Theoretical Biology 2025}.}
\talk{2024/09/12}{Proposal of morphological quantitative analysis methods using deep learning. \emph{Annual Meeting of the Japanese Society for Mathematical Biology 2024}.}
\talk{2024/03/28}{Quantification of mouse behavior under social defeat stress using mathematical models. \emph{CPSY TOKYO 2024}.}
\talk{2024/02/20}{Mathematical model of rumination in the brain using Bayesian inference. \emph{Spring School for Theoretical Biology 2024}.}
\talk{2024/01/06}{Mathematical model of rumination in the brain using free energy principle. \emph{11th Annual Meeting of the Quantitative Biology Society}.}
\talk{2023/12/16}{Unveiling the enigmatic Middle Devonian vertebrate, Palaeospondylus. \emph{The 3rd Asia Evo Conference Symposium}.}
\talk{2023/08/08}{Development of a quantitative analysis method for biological morphology using deep learning. \emph{Statphys28 Satellite meeting: Statistical Physics and Information-Processing in Living Systems}.}
\talk{2022/12/15}{Morphological feature extraction of primate mandibles using variational autoencoders. \emph{10th Annual Meeting of the Quantitative Biology Society}.}
\talk{2022/09/30}{A method for morphological feature extraction based on variational auto-encoder. \emph{60th Annual Meeting of the Biophysical Society of Japan}.}
\talk{2022/06/01}{A deep learning approach for the shape analysis of the primates mandible. \emph{Euro Evo Devo 2022}.}
\talk{2020/01/31}{Verification of classification methods for 2D OCT images using two-stage networks. \emph{2nd Japanese Society for Medical AI Academic Conference}.}

\section{Research Projects}
\begin{itemize}
  \item \textbf{Morpho-VAE}: Landmark-free morphological feature extraction with variational autoencoders.
  \item \textbf{mindGAP}: Population--individual gap in depressive symptom dynamics via energy landscape analysis.
  \item \textbf{KANADE}: Disentanglement of batch effects and biological signals in scRNA-seq.
  \item \textbf{DOLO}: Multi-individual Drosophila larva tracking with YOLO.
\end{itemize}

\vspace{0.35em}
{\footnotesize\color{muted}
Also online: \href{https://masa10223.github.io/en/publication/}{Publications} $\cdot$
\href{https://researchmap.jp/masato_tsutsumi?lang=en}{researchmap} $\cdot$
\href{https://scholar.google.co.jp/citations?user=sLCyYs8AAAAJ}{Google Scholar}
}

\end{document}
